\documentclass[11pt]{article}

\usepackage[margin=1in]{geometry}
\usepackage[T1]{fontenc}
\usepackage[utf8]{inputenc}
\usepackage{lmodern}
\usepackage{xcolor}
\usepackage{booktabs}
\usepackage{longtable}
\usepackage{array}
\usepackage{amsmath,amssymb}
\usepackage{enumitem}
\usepackage{hyperref}
\usepackage{pgfplotstable}
\usepackage{etoolbox}
\usepackage{shellesc}

\hypersetup{colorlinks=true, linkcolor=blue, urlcolor=blue}
\definecolor{accent}{RGB}{25,56,99}

\setlength{\parskip}{0.45em}
\setlength{\parindent}{0pt}

% ------------------------------------------------------------------
% USER SETTINGS
% ------------------------------------------------------------------
\newcommand{\ReportFreq}{0.01}
\newcommand{\CompareXLSX}{compare_cond_unorm__0p01Hz.xlsx}

% Temporary CSVs generated from workbook sheets
\newcommand{\TmpAll}{tmp_compare_ALL.csv}
\newcommand{\TmpSample}{tmp_cmp_SAMPLE.csv}
\newcommand{\TmpSite}{tmp_cmp_SITE.csv}
\newcommand{\TmpAllOnly}{tmp_cmp_ALL.csv}
\newcommand{\TmpSettings}{tmp_settings.csv}

% ------------------------------------------------------------------
% Convert XLSX sheets to CSV at compile time
% Requires: pdflatex --shell-escape
% Requires: xlsx2csv installed
% ------------------------------------------------------------------
\ShellEscape{xlsx2csv -n compare_ALL "\CompareXLSX" "\TmpAll"}
\ShellEscape{xlsx2csv -n cmp_SAMPLE "\CompareXLSX" "\TmpSample"}
\ShellEscape{xlsx2csv -n cmp_SITE "\CompareXLSX" "\TmpSite"}
\ShellEscape{xlsx2csv -n cmp_ALL "\CompareXLSX" "\TmpAllOnly"}
\ShellEscape{xlsx2csv -n settings "\CompareXLSX" "\TmpSettings"}

% ------------------------------------------------------------------
% PGFPLOTSTABLE SETTINGS
% ------------------------------------------------------------------
\pgfplotstableset{
  col sep=comma,
  string type,
  empty cells with={},
}

\pgfplotstableset{
  reportstyle/.style={
    every head row/.style={
      before row=\toprule,
      after row=\midrule
    },
    every last row/.style={
      after row=\bottomrule
    }
  }
}

\begin{document}

\begin{center}
{\LARGE\bfseries\color{accent}
Comparison of Power-Law and Logistic Models for Conductivity--Saturation Relationships}\\[0.4em]
{\large UNORM analysis of SIP column experiments at \ReportFreq\ Hz}
\end{center}

\section{Objective}

This report summarizes the conductivity-only UNORM comparison workflow using tables read directly from the workbook
\texttt{\CompareXLSX}.

The goal is to determine whether the conductivity--saturation relationship is better described by a power-law model or a logistic model.

The analysis compares the two candidate forms:
\[
\sigma = aS^b
\]
and
\[
\sigma = c + \frac{L}{1+\exp[-k(S-x_0)]}
\]

using the following model-selection statistics:

\begin{itemize}[leftmargin=1.5em]
\item residual standard deviation $\sigma$
\item Gaussian log-likelihood $\ln L$
\item Akaike Information Criterion (AIC)
\end{itemize}

The preferred model is the one with the smaller AIC.

\section{Input workbook}

This report is designed to compile in the same folder as:

\begin{itemize}[leftmargin=1.5em]
\item \texttt{\CompareXLSX}
\end{itemize}

The following sheets are read automatically from the workbook:

\begin{itemize}[leftmargin=1.5em]
\item \texttt{compare\_ALL}
\item \texttt{cmp\_SAMPLE}
\item \texttt{cmp\_SITE}
\item \texttt{cmp\_ALL}
\item \texttt{settings}
\end{itemize}

\section{Analysis structure}

The conductivity comparison is performed at three aggregation levels:

\begin{longtable}{p{0.20\textwidth} p{0.65\textwidth}}
\toprule
\textbf{Level} & \textbf{Description}\\
\midrule
\endhead
Sample & Each run is analyzed independently.\\
Site & All runs belonging to the same site are pooled together.\\
All & All available runs are combined into one dataset.\\
\bottomrule
\end{longtable}

Normalization, when enabled in the workflow, is applied only over the relevant group range.  
That means sample-level fits use sample-specific scaling, site-level fits use site-specific scaling, and all-data fits use the full combined range.

\section{Model comparison criteria}

For each fitted group, the workflow reports:

\begin{longtable}{p{0.24\textwidth} p{0.62\textwidth}}
\toprule
\textbf{Metric} & \textbf{Interpretation}\\
\midrule
\endhead
$R^2$ & Fraction of variance explained by the model.\\
SSE & Sum of squared residuals. Smaller values indicate a closer fit.\\
$\sigma$ & Residual standard deviation, estimated as $\sqrt{SSE/n}$.\\
$\ln L$ & Gaussian log-likelihood. Larger values indicate a better fit.\\
AIC & Information criterion balancing fit quality and model complexity. Smaller values are preferred.\\
\bottomrule
\end{longtable}

The power-law model has two fitted parameters, while the logistic model has four.  
AIC therefore penalizes the logistic model more strongly unless its improvement in fit is large enough to justify the extra flexibility.

\section{Workbook settings}

\pgfplotstabletypeset[
  reportstyle,
  columns={key,value},
  columns/key/.style={column name=Setting, string type, column type={p{0.28\textwidth}}},
  columns/value/.style={column name=Value, string type, column type={p{0.56\textwidth}}}
]{\TmpSettings}

\section{Automatic summary from workbook output}

The following tables are read directly from \texttt{\CompareXLSX}.  
This makes the report reproducible from the exported pipeline results.

\subsection{All fitted groups}

\pgfplotstabletypeset[
  reportstyle,
  columns={
    level,
    group,
    space,
    n_valid,
    winner,
    power_sigma_hat,
    power_lnL,
    power_aic,
    logistic_amp_L,
    logistic_sigma_hat,
    logistic_lnL,
    logistic_aic,
    delta_aic_logistic_minus_power
  },
  columns/level/.style={column name=Level, string type, column type={p{0.10\textwidth}}},
  columns/group/.style={column name=Group, string type, column type={p{0.15\textwidth}}},
  columns/space/.style={column name=Space, string type, column type={p{0.08\textwidth}}},
  columns/n_valid/.style={column name=$n$, column type={r}},
  columns/winner/.style={column name=Winner, string type, column type={p{0.10\textwidth}}},
  columns/power_sigma_hat/.style={column name=Power $\sigma$, fixed, precision=4, column type={r}},
  columns/power_lnL/.style={column name=Power $\ln L$, fixed, precision=3, column type={r}},
  columns/power_aic/.style={column name=Power AIC, fixed, precision=3, column type={r}},
  columns/logistic_amp_L/.style={column name=Logistic $L$, fixed, precision=4, column type={r}},
  columns/logistic_sigma_hat/.style={column name=Logistic $\sigma$, fixed, precision=4, column type={r}},
  columns/logistic_lnL/.style={column name=Logistic $\ln L$, fixed, precision=3, column type={r}},
  columns/logistic_aic/.style={column name=Logistic AIC, fixed, precision=3, column type={r}},
  columns/delta_aic_logistic_minus_power/.style={column name=$\Delta$AIC, fixed, precision=3, column type={r}}
]{\TmpAll}

\section{Condensed interpretation guide}

The last column in the table above is
\[
\Delta AIC = AIC_{\text{logistic}} - AIC_{\text{power}}
\]

Interpretation:

\begin{itemize}[leftmargin=1.5em]
\item $\Delta AIC < 0$ means the logistic model is preferred.
\item $\Delta AIC > 0$ means the power-law model is preferred.
\item Larger negative values indicate stronger support for the logistic model.
\item Larger positive values indicate stronger support for the power-law model.
\end{itemize}

\section{Sample-level results}

\pgfplotstabletypeset[
  reportstyle,
  columns={
    group,
    n_valid,
    winner,
    power_aic,
    logistic_aic,
    delta_aic_logistic_minus_power,
    power_equation,
    logistic_equation
  },
  columns/group/.style={column name=Sample, string type, column type={p{0.14\textwidth}}},
  columns/n_valid/.style={column name=$n$, column type={r}},
  columns/winner/.style={column name=Winner, string type, column type={p{0.10\textwidth}}},
  columns/power_aic/.style={column name=Power AIC, fixed, precision=3, column type={r}},
  columns/logistic_aic/.style={column name=Logistic AIC, fixed, precision=3, column type={r}},
  columns/delta_aic_logistic_minus_power/.style={column name=$\Delta$AIC, fixed, precision=3, column type={r}},
  columns/power_equation/.style={column name=Power equation, string type, column type={p{0.20\textwidth}}},
  columns/logistic_equation/.style={column name=Logistic equation, string type, column type={p{0.24\textwidth}}}
]{\TmpSample}

\section{Site-level results}

\pgfplotstabletypeset[
  reportstyle,
  columns={
    group,
    n_valid,
    winner,
    power_aic,
    logistic_aic,
    delta_aic_logistic_minus_power,
    power_equation,
    logistic_equation
  },
  columns/group/.style={column name=Site, string type, column type={p{0.14\textwidth}}},
  columns/n_valid/.style={column name=$n$, column type={r}},
  columns/winner/.style={column name=Winner, string type, column type={p{0.10\textwidth}}},
  columns/power_aic/.style={column name=Power AIC, fixed, precision=3, column type={r}},
  columns/logistic_aic/.style={column name=Logistic AIC, fixed, precision=3, column type={r}},
  columns/delta_aic_logistic_minus_power/.style={column name=$\Delta$AIC, fixed, precision=3, column type={r}},
  columns/power_equation/.style={column name=Power equation, string type, column type={p{0.20\textwidth}}},
  columns/logistic_equation/.style={column name=Logistic equation, string type, column type={p{0.24\textwidth}}}
]{\TmpSite}

\section{All-data result}

\pgfplotstabletypeset[
  reportstyle,
  columns={
    group,
    space,
    n_valid,
    winner,
    power_sigma_hat,
    power_lnL,
    power_aic,
    logistic_amp_L,
    logistic_sigma_hat,
    logistic_lnL,
    logistic_aic,
    delta_aic_logistic_minus_power,
    power_equation,
    logistic_equation
  },
  columns/group/.style={column name=Group, string type, column type={p{0.10\textwidth}}},
  columns/space/.style={column name=Space, string type, column type={p{0.08\textwidth}}},
  columns/n_valid/.style={column name=$n$, column type={r}},
  columns/winner/.style={column name=Winner, string type, column type={p{0.10\textwidth}}},
  columns/power_sigma_hat/.style={column name=Power $\sigma$, fixed, precision=4, column type={r}},
  columns/power_lnL/.style={column name=Power $\ln L$, fixed, precision=3, column type={r}},
  columns/power_aic/.style={column name=Power AIC, fixed, precision=3, column type={r}},
  columns/logistic_amp_L/.style={column name=Logistic $L$, fixed, precision=4, column type={r}},
  columns/logistic_sigma_hat/.style={column name=Logistic $\sigma$, fixed, precision=4, column type={r}},
  columns/logistic_lnL/.style={column name=Logistic $\ln L$, fixed, precision=3, column type={r}},
  columns/logistic_aic/.style={column name=Logistic AIC, fixed, precision=3, column type={r}},
  columns/delta_aic_logistic_minus_power/.style={column name=$\Delta$AIC, fixed, precision=3, column type={r}},
  columns/power_equation/.style={column name=Power equation, string type, column type={p{0.18\textwidth}}},
  columns/logistic_equation/.style={column name=Logistic equation, string type, column type={p{0.22\textwidth}}}
]{\TmpAllOnly}

\section{Physical interpretation}

Conductivity in partially saturated porous media may follow either of two broad patterns.

\subsection{Power-law behavior}

A power-law fit suggests a smooth monotonic scaling between saturation and conductivity, without a distinct transition zone.

\subsection{Logistic behavior}

A logistic fit suggests threshold-like behavior. This may reflect:

\begin{itemize}[leftmargin=1.5em]
\item formation of connected water pathways
\item onset of pore-network connectivity
\item rapid rise in conductivity over a limited saturation interval
\item stabilization toward an upper plateau at high saturation
\end{itemize}

In many partially saturated systems, conductivity is more sensitive than resistivity to such connectivity transitions.

\section{How to read the exported plots}

The comparison workflow also writes plots for each fitted group.  
Each plot contains:

\begin{itemize}[leftmargin=1.5em]
\item measured conductivity data
\item fitted power-law curve
\item fitted logistic curve
\item optional labels with $\sigma$, $\ln L$, and AIC
\end{itemize}

These plots should be interpreted together with the AIC table.  
A visually flexible curve is not automatically preferred unless its AIC is smaller.

\section{Conclusion}

This report provides a reproducible summary of the UNORM conductivity comparison workflow.

The key decision rule is:

\[
\text{Choose the model with the smaller AIC}
\]

while also reviewing:

\begin{itemize}[leftmargin=1.5em]
\item residual standard deviation
\item Gaussian log-likelihood
\item fitted equations
\item plotted curve shapes
\end{itemize}

Together, these outputs help determine whether conductivity is better represented by smooth power-law scaling or by threshold-driven logistic behavior.

\end{document}
